\section{CONCLUSIONES Y TRABAJOS FUTUROS}

En este trabajo se presentó un modelo computacional de aprendizaje adaptativo para la enseñanza de teoría musical, el cual integra un Grafo de Conocimiento Dirigido Acíclico (DAG) para representar explícitamente las dependencias pedagógicas del dominio y un agente basado en \textit{Deep Reinforcement Learning} para recomendar dinámicamente actividades de aprendizaje de acuerdo con el progreso individual del estudiante. El modelo fue implementado mediante la plataforma distribuida \textbf{SheetMusic}, permitiendo materializar la propuesta en un entorno real de apoyo al proceso de enseñanza y aprendizaje.

Los resultados de la validación experimental muestran que el agente de recomendación fue capaz de adaptar la secuencia y dificultad de los ejercicios según el desempeño de cada estudiante, favoreciendo una progresión sostenida del nivel de conocimiento y manteniendo una participación continua durante el período de evaluación. Asimismo, la evaluación de usabilidad y aceptación tecnológica evidenció una percepción favorable de la plataforma por parte de sus usuarios, respaldando la viabilidad de la propuesta para su utilización en contextos educativos.

Desde la perspectiva de la Ingeniería Informática, uno de los principales aportes del trabajo corresponde a la integración de una representación explícita del conocimiento pedagógico mediante un DAG con un modelo de aprendizaje por refuerzo profundo, desacoplando la estructura del dominio, el proceso de toma de decisiones y la implementación de la plataforma. Esta arquitectura facilita la evolución independiente de cada componente, permitiendo incorporar nuevos contenidos, estrategias de recomendación o modelos de inteligencia artificial sin modificar la infraestructura principal de \textbf{SheetMusic}.

Con el propósito de favorecer la reproducibilidad científica y la continuidad del proyecto, el código fuente de la plataforma, del motor de recomendación y de los microservicios desarrollados se encuentra disponible públicamente en la organización oficial de GitHub: \url{https://github.com/Sheetmusic-Team}.

Como trabajo futuro se contempla ampliar el Grafo de Conocimiento Musical para cubrir la totalidad de las unidades curriculares de teoría musical, incorporar nuevos modelos de aprendizaje por refuerzo y comparar su desempeño con PPO, integrar evaluación instrumental en tiempo real mediante procesamiento digital de señales (DSP) y validar la plataforma con una muestra significativamente mayor de estudiantes pertenecientes a distintos establecimientos educacionales del país.